% sections/abstract.tex — Abstract
% v1 from notes/paper_abstract_draft.md, integrated 2026-05-13

Vision-Language-Action (VLA) models achieve up to 97\% task success rate on simulation manipulation benchmarks, with paraphrase invariance interpreted as evidence of language understanding. We question this interpretation. Through a three-layer signal-flow diagnostic on a strong VLA backbone trained on LIBERO-Spatial, we find that backbone hidden states correlate strongly with instruction embeddings ($r = +0.76$, $p < 0.001$) while demonstration trajectories correlate only weakly with the same embeddings ($r = -0.27$, $p = 0.96$) --- an approximately $8\times$ squared-correlation gap localizing language information loss to the demonstration-supervision stage rather than to backbone capacity. We formalize this via Theorem 3$'$, an architecture-agnostic upper bound: any demonstration-supervised policy minimizing a distribution-matching loss is bounded by $I(\Astar; \ell \mid o) \le I_{\dataD}(a; \ell \mid o)$. The bound applies uniformly across $L_1$/$L_2$ regression, Conditional Flow Matching, Shortcut Flow Matching, and discretized cross-entropy heads. OpenVLA-OFT's 97\% SR is consistent with two scenarios within the bound --- low-magnitude language preservation or entirely visual grounding --- neither distinguishable by benchmark SR alone. We document a Conditional Flow Matching decoder sample-quality collapse failure mode (7/7 action dimensions fail under language-orthogonal supervision, oracle ablation rules out moving-target hypothesis), characterizing how decoder family choice determines failure mode within the bound. Our contribution is a diagnostic methodology and theoretical bound; we explicitly do not propose an architectural fix, as bypassing the bound requires interventions outside the demonstration-supervised loss.


